% 1. Describe the overall context + bring several background references
% 2. Show the research gap (why to do this particular research)
% 3. Bring up your research question(s)
% 4. Describe the structure of the work + mention research method

% Research methods: literature review, design science prototype, action-validity evaluation

\todo{A paragraph filled with ancient web development and the modern LLM revolution, leads to topic}

\subsection{Why this topic}

\todo{Find supporting materials for this and make each bullet a point}
\begin{itemize}
    \item Cognitive load and inefficiency of pure conversational interfaces in complex, state-dependent web tasks.
    \item Current Agentic/Generative UI systems often have to infer available actions from unstructured context, human-facing pages, or weakly structured APIs.
    \item Utilizing HTML hypermedia representations that expose state-dependent affordances as a trusted information layer for state-aware agentic UI generation.
\end{itemize}

%\cite{jarvi_algorithms_2009}
%\cite{yang2025agentic}
%\citet{jarvi_algorithms_2009}

\subsection{Research Questions}

\begin{enumerate}
    \item[] \textbf{RQ1:} What are the fundamental limitations of CUIs that necessitate Agentic UIs, and how do state-of-the-art Generative UI frameworks differ in terms of their benefits and drawbacks?
    
    % \item[] \textbf{RQ2:} How can a system architecture be designed to parse backend hypermedia constraints (application state and affordances) and deterministically map them to frontend components based on user intent?
    \item[] \textbf{RQ2:}  How can HTML hypermedia affordances, such as links, forms, actions, and element states, be extracted and represented as a trusted constraint layer for state-aware agentic UI generation?
    

    % \item[] \textbf{RQ3:} Measured via computational HCI models (e.g., Keystroke-Level Model), to what extent do hypermedia-generated GUIs reduce interaction complexity and operation step counts for complex web tasks compared to standard conversational interfaces?
    \item[] \textbf{RQ3:} In a simulated state-dependent workflow, to what extent does an HTML hypermedia-constrained agentic UI reduce invalid action suggestions compared with an unconstrained or weakly structured baseline?
\end{enumerate}

\todo{Point out the contribution of this work: a framework comparison, an HTML hypermedia-constrained architecture, and a prototype evaluation focused on action validity rather than general UI quality.}


\subsection{Structure of the work}

\todo{Describe the structure of the work and tie research questions to it. 

Better do this at the final stage}
